% !TEX program = xelatex
% !TEX output = 回归分析前三章.pdf
\PassOptionsToPackage{fontset=fandol}{ctex}
\documentclass[lang=cn,10pt]{elegantbook}
\usepackage{dirtree}
\usepackage{needspace}

\title{回归分析}
\subtitle{前三章学习笔记}
\author{东邪}
\extrainfo{不要以为抹消过去，重新来过，即可发生什么改变。—— 比企谷八幡}
\cover{images/cover.png}
\logo{images/logo-blue.png}
\date{2026年10月6日}
\version{1.0}
\bioinfo{内容}{定义、定理与必要公式}
\setcounter{tocdepth}{1}

% 数理统计封面原图为 1281 x 1024；固定同高的图片区，保持本图比例。
% 仅在主文件补足画布留白，沿用原类文件的横带、文字、标志和座右铭布局。
\makeatletter
\patchcmd{\maketitle}
  {\includegraphics[width=\linewidth]{\@cover}}
  {\begin{minipage}[c][0.8\linewidth][c]{\linewidth}
     \centering\includegraphics[width=\linewidth]{\@cover}
   \end{minipage}}
  {}{\ClassError{elegantbook}{Cover layout patch failed}{Check the cover command.}}
% 将同款标志直接排在右侧信息栏，避免绝对定位受编译器缓存影响。
\patchcmd{\maketitle}
  {\begin{tikzpicture}[remember picture,overlay]
   \begin{pgfonlayer}{background}
   \node[opacity=0.8,
         anchor=south east,
         outer sep=0pt,
         inner sep=0pt] at ($(current page.south east) +(-0.8in,1.5in)$) {
   \ifcsname @logo\endcsname\includegraphics[width=4.2cm]{\@logo}\fi};
   \end{pgfonlayer}
   \end{tikzpicture}}
  {\centering\includegraphics[width=4.2cm]{\@logo}}
  {}{\ClassError{elegantbook}{Cover logo patch failed}{Check the logo command.}}
\makeatother

\usepackage{mathtools}
\newcommand{\R}{\mathbb{R}}
\newcommand{\E}{\mathbb{E}}
\DeclareMathOperator{\Var}{Var}
\DeclareMathOperator{\Cov}{Cov}
\DeclareMathOperator{\rank}{rank}
\DeclareMathOperator{\tr}{tr}
\DeclareMathOperator*{\argmin}{arg\,min}
\newcommand{\XX}{\mathbf{X}}
\newcommand{\bb}{\widehat{\beta}}
\newcommand{\uh}{\widehat{u}}
\hypersetup{pdftitle={回归分析：前三章学习笔记},pdfauthor={东邪}}

\begin{document}
\maketitle
\frontmatter
\tableofcontents
\mainmatter
\chapter{什么是回归}

\section*{回归分析的功能结构}
\phantomsection
\addcontentsline{toc}{section}{回归分析的功能结构}
\begingroup\small
\dirtree{%
.1 回归分析：用解释变量研究响应变量的条件分布.
.2 研究对象：随机向量 $(Y,X')'$ 与观测样本.
.3 条件均值回答平均怎样变化；条件方差回答波动怎样变化.
.2 描述关系：哪些变量共同变化，怎样组织模型？.
.3 工具：协方差、相关系数、领域理论、DAG（第1章）.
.3 区分相关、预测与因果；因果解释还需识别假设.
.2 预测响应：给定 $X$，怎样预测 $Y$？.
.3 全部平方可积函数中选 CEF，最小化 MSE（第2章）.
.3 限制为仿射函数，得到最优线性投影（第2章）.
.3 用样本矩替代总体矩，得到 OLS（第3章）.
.2 量化不确定性：系数可靠吗，限制成立吗？.
.3 工具：外生性、球形方差、Gauss--Markov（第3章）.
.3 正态假设给出精确分布、$t/F$ 检验与区间（第3章）.
.3 已知线性限制可用 CLS，提高估计效率（第3章）.
.3 大样本理论与高维方法（第4、5章，本次正文未展开）.
.2 检查模型：所拟合的对象是否回答原问题？.
.3 正确设定：CEF 与线性投影何时一致（第2章）.
.3 检查秩、外生性、同方差、正态性各自支持什么结论.
.3 拟合优度描述样本解释率；预测需关注样本外表现.
.2 连接各部分的工具链.
.3 条件期望与矩：把预测问题写成正交条件.
.3 线性代数与矩阵求导：从正交条件到投影和 OLS.
.3 正态向量与二次型：从投影到抽样分布和检验.
}
\endgroup
\smallskip
这棵树是按功能重组的个人学习框架。划分参考 CMU 的
\href{https://www.stat.cmu.edu/~cshalizi/mreg/15/}{Cosma Shalizi：Modern Regression（2015）}
中预测、估计、推断、假设检查与模型选择的安排，以及
\href{https://www.stat.cmu.edu/~larry/=stat401/}{Larry Wasserman：Modern Regression（2017）}
的课程进路；上树的中文分组和本讲义章次映射由东邪整理。
\par\smallskip
前三章把“研究关系”连到“选择总体预测器”，再连到“样本估计及推断”。预测损失、模型设定与因果识别是不同问题，不能仅凭显著系数或较高 $R^2$ 替代后两者的论证。
\clearpage
\begin{introduction}[本章的作用]
\item 用相关系数度量线性共同变动。
\item 用条件分布与领域知识提出模型。
\item 用 DAG 检查路径与控制变量的作用。
\end{introduction}
\section{回归对象与条件均值}

% SOURCE: 原文未编号概念或必要公式; page-0005.tex
\begin{definition}[回归分析与 CEF]\label{reg:entry:1}
响应变量 $Y$ 为标量，解释变量 $X$ 为向量。回归分析研究 $Y$ 随 $X$ 的取值而变化的规律，尤其是条件分布及条件均值。若 $\E|Y|<\infty$，记
\[\mu(x)=\E[Y\mid X=x],\qquad Y=\mu(X)+e,\qquad \E[e\mid X]=0.\]
在连续情形且 $f_X(x)>0$ 时，$f_{Y\mid X}(y\mid x)=f_{X,Y}(x,y)/f_X(x)$，
\[\mu(x)=\int_{\R}y f_{Y\mid X}(y\mid x)\,dy.\]
可用于描述关系和预测；把关系解释为因果效应，需要额外识别条件。
\end{definition}

% SOURCE: 原文未编号概念或必要公式; page-0007.tex 至 page-0010.tex
\begin{proposition}[二元正态的回归关系]\label{reg:entry:2}
设 $(X,Y)$ 联合正态，均值为 $\mu_X,\mu_Y$，方差为 $\sigma_X^2,\sigma_Y^2>0$，相关参数 $|\rho|<1$。令 $z_X=(x-\mu_X)/\sigma_X$、$z_Y=(y-\mu_Y)/\sigma_Y$，其联合密度为
\[f(x,y)=\frac{\exp\!\left[-\frac{z_X^2-2\rho z_Xz_Y+z_Y^2}{2(1-\rho^2)}\right]}{2\pi\sigma_X\sigma_Y\sqrt{1-\rho^2}}.\]
边缘分布分别为 $N(\mu_X,\sigma_X^2)$、$N(\mu_Y,\sigma_Y^2)$，并且
\begin{align*}
Y\mid X&\sim N\!\left(\mu_Y+\rho\frac{\sigma_Y}{\sigma_X}(X-\mu_X),\,\sigma_Y^2(1-\rho^2)\right),\\
X\mid Y&\sim N\!\left(\mu_X+\rho\frac{\sigma_X}{\sigma_Y}(Y-\mu_Y),\,\sigma_X^2(1-\rho^2)\right).
\end{align*}
两个方向的斜率为 $\beta_{Y\mid X}=\rho\sigma_Y/\sigma_X$ 与 $\beta_{X\mid Y}=\rho\sigma_X/\sigma_Y$；乘积为 $\rho^2$，当两方差相等时两斜率相同。
\end{proposition}

\section{相关性}

% SOURCE: 定义 1.2.1; page-0010.tex
\begin{definition}[相关系数]\label{reg:entry:3}
若 $X,Y$ 的二阶矩有限且方差均正，定义
\[\rho_{XY}=\frac{\Cov(X,Y)}{\sqrt{\Var(X)\Var(Y)}}
=\frac{\E[XY]-\E[X]\E[Y]}{\sqrt{\Var(X)\Var(Y)}}.\]
其中 $\Cov(X,Y)=\E[(X-\E X)(Y-\E Y)]$。
\end{definition}

% SOURCE: 定理 1.2.1; page-0011.tex
\begin{theorem}[相关系数与 Cauchy--Schwarz]\label{reg:entry:4}
设 $X,Y$ 二阶矩有限。
\begin{enumerate}
\item 若 $X,Y$ 独立，则协方差为零；方差正时 $\rho_{XY}=0$。一般逆命题不成立。
\item $\E|XY|\leq\sqrt{\E[X^2]\E[Y^2]}$。
\item 若两方差均正，则 $|\rho_{XY}|\leq1$。
\end{enumerate}
联合正态时，零相关等价于独立。
\end{theorem}

% SOURCE: 原文未编号概念或必要公式; page-0013.tex
\begin{definition}[样本相关系数]\label{reg:entry:5}
对于配对样本 $(X_i,Y_i)_{i=1}^n$，令 $\bar X=n^{-1}\sum_i X_i$、$\bar Y=n^{-1}\sum_iY_i$，$s_X^2=n^{-1}\sum_i(X_i-\bar X)^2$、$s_Y^2=n^{-1}\sum_i(Y_i-\bar Y)^2$。若 $s_Xs_Y>0$，定义
\[R=\frac{n^{-1}\sum_i X_iY_i-\bar X\bar Y}{s_Xs_Y}
=\frac{\sum_i(X_i-\bar X)(Y_i-\bar Y)}{\sqrt{\sum_i(X_i-\bar X)^2\sum_i(Y_i-\bar Y)^2}}.\]
\end{definition}

% SOURCE: 定理 1.2.2; page-0012.tex
\begin{theorem}[样本相关系数的一致性]\label{reg:entry:6}
若配对样本独立同分布，$\E X^2,\E Y^2<\infty$ 且总体两方差均正，则在样本相关系数定义的事件上，$R\xrightarrow{P}\rho_{XY}$；该事件的概率趋于 $1$。
\end{theorem}

% SOURCE: 原文未编号概念或必要公式; page-0014.tex
\begin{proposition}[相关性与因果关系]\label{reg:entry:7}
相关系数只衡量线性关联。零相关不能排除非线性依赖；相关不等于因果，共同原因、选择机制等都可能产生关联。用回归系数解释因果效应时，必须另行说明研究设计和识别假设。
\end{proposition}

\section{构建回归模型的步骤}

% SOURCE: 原文未编号概念或必要公式; page-0014.tex 至 page-0016.tex
\begin{proposition}[四步建模及两类假设]\label{reg:entry:8}
\begin{enumerate}
\item 收集数据并作描述性分析，寻找需要解释的模式。
\item 结合领域知识建立理论模型，确定对象、变量和作用机制。
\item 选择函数形式、未知参数及误差条件，进行估计和检验。
\item 用模型解释或预测，并检查结论是否适用于原问题。
\end{enumerate}
要分清函数形式及变量选择等\textbf{模型设定假设}，与估计和推断所需的\textbf{统计假设}；估计量表现良好并不能证明函数形式正确。
\end{proposition}

\section{有向无环图与路径}

% SOURCE: 原文未编号概念或必要公式; page-0016.tex
\begin{definition}[DAG]\label{reg:entry:9}
有向无环图（directed acyclic graph）以节点表示变量，以有向边表示假定的因果方向，且不存在有向循环。讲义用实心节点表示可观测变量、空心节点表示不可观测变量。箭头编码的是模型假设，不能仅从相关系数推出。
\end{definition}

% SOURCE: 原文未编号概念或必要公式; page-0016.tex 至 page-0018.tex
\begin{proposition}[三种基本路径的功能]\label{reg:entry:10}
链 $X\to M\to Y$ 表示中介路径；分叉 $X\leftarrow C\to Y$ 表示共同原因路径；碰撞结构 $X\to C\leftarrow Y$ 的中间节点称为碰撞点。
按 DAG 的路径阻断规则，前两条路径在未控制中间节点时开放，控制中间节点则阻断；碰撞路径默认阻断，控制碰撞点或其后代可能使它开放。
在采用图的 Markov 条件时，所有路径被阻断（d-separation）才给出相应条件独立；一条路径开放只表示图不强制独立，不保证实际相关系数非零。控制混杂变量可阻断相应后门路径，但能否识别因果效应须检查全部相关路径和识别条件。
\end{proposition}


\chapter{条件期望函数}

\begin{introduction}[本章的作用]
\item 条件期望是平方损失下的最优预测。
\item 线性投影把全部函数的优化限制为仿射函数。
\item 正确设定连接投影系数与真实条件均值参数。
\end{introduction}
本章 $Y\in\R$，总体解释向量 $X\in\R^K$；含截距时统一写成 $X=(1,X_1,\ldots,X_k)'$，$K=k+1$。
\section{条件期望、重期望与最优预测}

% SOURCE: 定义 2.1.1; page-0019.tex
\begin{definition}[条件期望函数]\label{reg:entry:11}
若 $\E|Y|<\infty$，定义 $\mu(x)=\E[Y\mid X=x]$，称为条件期望函数（CEF）；作为随机变量写成 $\mu(X)=\E[Y\mid X]$。连续情形可用条件密度积分，离散情形用条件分布求和。等式按几乎处处意义理解。
\end{definition}

% SOURCE: 原文未编号概念或必要公式; page-0020.tex
\begin{proposition}[条件期望的基本性质]\label{reg:entry:12}
在所涉及随机变量可积的条件下，对于常数 $a,b$ 和可测函数 $g$，
\[\E[aY+b\mid X]=a\E[Y\mid X]+b,\qquad \E[g(X)Y\mid X]=g(X)\E[Y\mid X].\]
若 $X,Y$ 独立，则 $\E[Y\mid X]=\E Y$，以及 $\E[g(X)Y\mid X]=g(X)\E Y$。乘积公式要求 $g(X)Y$ 可积。
\end{proposition}

% SOURCE: 定义 2.1.2; page-0020.tex
\begin{definition}[均值独立]\label{reg:entry:13}
若 $\E[Y\mid X]=\E Y$，称 $Y$ 均值独立于 $X$。此性质通常不对称。相互独立推出两个方向的均值独立；若 $X,Y$ 二阶矩有限，任一方向的均值独立均推出 $\Cov(X,Y)=0$。这些逆命题一般不成立。
\end{definition}

% SOURCE: 定理 2.1.1; page-0021.tex
\begin{theorem}[重期望公式 I]\label{reg:entry:14}
若 $\E|Y|<\infty$，则对任意随机向量 $X$，
\[\E_X\bigl[\E_Y[Y\mid X]\bigr]=\E_Y[Y].\]
\end{theorem}

% SOURCE: 定理 2.1.2; page-0022.tex
\begin{theorem}[重期望公式 II]\label{reg:entry:15}
对任意可测函数 $G(X,Y)$，若 $\E|G(X,Y)|<\infty$，则
\[\E_X\bigl[\E_Y[G(X,Y)\mid X]\bigr]=\E_{X,Y}[G(X,Y)].\]
\end{theorem}

% SOURCE: 定义 2.1.3; page-0022.tex
\begin{definition}[均方误差]\label{reg:entry:16}
用可测函数 $g(X)$ 预测 $Y$，定义均方误差
\[\operatorname{MSE}(g)=\E[(Y-g(X))^2].\]
以下最优预测结论要求 $Y$ 及候选预测函数平方可积。
\end{definition}

% SOURCE: 定理 2.1.3; page-0022.tex
\begin{theorem}[最优预测]\label{reg:entry:17}
设 $\E Y^2<\infty$，$\mathbb F=\{g:\R^K\to\R:\ g\text{可测},\ \E[g(X)^2]<\infty\}$。则
\[\mu(X)=\E[Y\mid X]=\argmin_{g\in\mathbb F}\operatorname{MSE}(g).\]
更准确地，最优预测在几乎处处意义下唯一，且
\[\operatorname{MSE}(g)=\E[(Y-\mu(X))^2]+\E[(\mu(X)-g(X))^2].\]
\end{theorem}

\section{条件方差与 CEF 误差}

% SOURCE: 定义 2.1.4; page-0023.tex
\begin{definition}[条件方差]\label{reg:entry:18}
若 $\E Y^2<\infty$，定义
\begin{align*}
\sigma^2(X)&=\Var(Y\mid X)=\E[(Y-\E[Y\mid X])^2\mid X]\\
&=\E[Y^2\mid X]-\bigl(\E[Y\mid X]\bigr)^2.
\end{align*}
在 $X=x$ 时记为 $\sigma^2(x)$。
\end{definition}

% SOURCE: 定理 2.1.4; page-0024.tex
\begin{theorem}[条件方差分解]\label{reg:entry:19}
若 $\E Y^2<\infty$，则
\[\Var(Y)=\E[\Var(Y\mid X)]+\Var(\E[Y\mid X]).\]
\end{theorem}

% SOURCE: 原文未编号概念或必要公式; page-0024.tex 至 page-0025.tex
\begin{definition}[CEF 误差]\label{reg:entry:20}
定义 $e=Y-\mu(X)$，其中 $\mu(X)=\E[Y\mid X]$。由此得 $Y=\mu(X)+e$。
反过来，若 $Y=\mu(X)+e$，$Y,e$ 可积、$\mu$ 可测且 $\E[e\mid X]=0$，则 $\mu(X)=\E[Y\mid X]$。
\end{definition}

% SOURCE: 定理 2.1.5; page-0024.tex
\begin{theorem}[CEF 误差的正交性]\label{reg:entry:21}
若 $\E|Y|<\infty$，$e=Y-\E[Y\mid X]$，则
\[\E[e\mid X]=0,\qquad \E e=0.\]
对任意可测函数 $h$，若 $\E|h(X)e|<\infty$，还有 $\E[h(X)e]=0$。乘积可积条件不可省略。
\end{theorem}

% SOURCE: 定义 2.1.5; page-0025.tex
\begin{definition}[同方差与异方差]\label{reg:entry:22}
对于平方可积的 CEF 误差，$\Var(e\mid X)=\sigma^2(X)$。若 $\sigma^2(x)=\sigma^2$ 不依赖 $x$（几乎处处），称误差同方差；若依赖于 $x$，称异方差。零条件均值并不推出同方差或统计独立。
\end{definition}

\section{线性 CEF 与最优线性预测}

% SOURCE: 定义 2.2.1; page-0026.tex
\begin{definition}[仿射函数]\label{reg:entry:23}
令 $X=(1,X_1,\ldots,X_k)'\in\R^K$，$K=k+1$，$\beta=(\beta_0,\ldots,\beta_k)'$。仿射函数类为
\[\mathbb A=\left\{g(X)=X'\beta=\beta_0+\sum_{j=1}^k\beta_jX_j:\ \beta\in\R^K\right\}.\]
把常数 $1$ 放入解释向量后，可统一称为对扩充向量 $X$ 的线性函数。
\end{definition}

% SOURCE: 定义 2.2.2; page-0026.tex
\begin{definition}[线性模型与线性 CEF]\label{reg:entry:24}
当 $g\in\mathbb A$ 时，$Y=g(X)+e=X'\beta+e$ 称为 $Y$ 对 $X$ 的线性模型。
若另有 $\E[e\mid X]=0$，则 $\E[Y\mid X]=X'\beta$，称为线性 CEF 模型。仅有代数分解并不保证条件均值线性。
\end{definition}

% SOURCE: 假设 2.2.1; page-0027.tex
\begin{proposition}[线性预测的矩与识别条件（假设 2.2.1）]\label{reg:entry:25}
设 $X\in\R^K$、$Y\in\R$，满足
\[\E Y^2<\infty,\qquad \E\|X\|^2<\infty,\qquad Q_{XX}:=\E[XX']\succ0.\]
前两条保证二阶矩和混合矩有限；正定性保证 $Q_{XX}$ 可逆，排除总体上的精确线性依赖。
\end{proposition}

% SOURCE: 定义 2.2.3; page-0028.tex
\begin{definition}[最优线性预测与投影系数]\label{reg:entry:26}
在上述矩条件下，最优线性预测（线性投影）定义为
\[L(Y\mid X)=X'\beta^*,\qquad \beta^*=\argmin_{\beta\in\R^K}\E[(Y-X'\beta)^2].\]
称 $\beta^*$ 为线性投影系数。目标函数为
\[S(\beta)=\E Y^2-2\beta'\E[XY]+\beta'Q_{XX}\beta.\]
\end{definition}

% SOURCE: 定理 2.2.1; page-0028.tex
\begin{theorem}[投影系数的存在唯一性]\label{reg:entry:27}
在假设 2.2.1 的有限二阶矩与正定性条件下，$\beta^*$ 存在且唯一，
\[\beta^*=Q_{XX}^{-1}\E[XY].\]
正规方程为 $Q_{XX}\beta^*=\E[XY]$，Hessian 为 $2Q_{XX}\succ0$。
\end{theorem}

% SOURCE: 定理 2.2.2; page-0029.tex
\begin{theorem}[投影误差的刻画]\label{reg:entry:28}
在假设 2.2.1 的条件下，令 $Y=X'\beta+e_L$，则
\[\beta=\beta^*\quad\Longleftrightarrow\quad\E[Xe_L]=0.\]
若 $X$ 含常数项，则 $\E e_L=0$。一般 $\E[e_L\mid X]$ 未必为零；CEF 误差与投影误差只在线性 CEF 的情形一致。
\end{theorem}

% SOURCE: 原文未编号概念或必要公式; page-0030.tex
\begin{proposition}[CEF 与线性投影的损失比较]\label{reg:entry:29}
设上述条件成立，并令 $e_L=Y-X'\beta^*$。则
\[\E[(Y-X'\beta^*)^2]=\E[(Y-\mu(X))^2]+\E[(\mu(X)-X'\beta^*)^2].\]
因此最优线性预测的 MSE 不小于 CEF 的 MSE；等号当且仅当 $\mu(X)=X'\beta^*$ 几乎处处。不能把两个预测器本身写成大小不等式。
\end{proposition}

\section{条件均值模型的正确设定}

% SOURCE: 定义 2.3.1; page-0031.tex
\begin{definition}[正确设定与错误设定]\label{reg:entry:30}
对 $Y\in\R$、$X\in\R^K$，若存在 $\beta^o\in\R^K$ 使
\[\E[Y\mid X]=X'\beta^o\quad\text{几乎处处},\]
称线性模型对条件均值正确设定；若不存在这样的参数，称错误设定。正确设定等价于存在 $\beta^o$ 使误差 $Y-X'\beta^o$ 的条件均值为零。
\end{definition}

% SOURCE: 定理 2.3.1; page-0031.tex
\begin{theorem}[正确设定连接 CEF 与投影]\label{reg:entry:31}
在假设 2.2.1 的矩条件下，若线性模型对 $\E[Y\mid X]$ 正确设定，则：
\begin{enumerate}
\item 存在 $\beta^o$ 和 $e=Y-X'\beta^o$，使 $Y=X'\beta^o+e$ 且 $\E[e\mid X]=0$。
\item 线性投影系数满足 $\beta^*=\beta^o$。
\end{enumerate}
对于连续解释变量 $X_j$，线性 CEF 给出偏导 $\partial\E[Y\mid X]/\partial X_j=\beta_j^o$；截距不属于这一偏导，因果解释仍需额外条件。
\end{theorem}

\section{拾遗：矩阵求导}

% SOURCE: 原文未编号概念或必要公式; page-0033.tex
\begin{definition}[梯度与导数排列]\label{reg:entry:32}
对可微标量函数 $g:\R^K\to\R$，采用列梯度
\[\frac{\partial g(x)}{\partial x}=(\partial g/\partial x_1,\ldots,\partial g/\partial x_K)'.\]
对列向量函数 $F$，Jacobian 的 $(i,j)$ 元为 $\partial F_i/\partial x_j$。行向量函数的按列梯度堆叠与此互为转置。
\end{definition}

% SOURCE: 定理 2.5.1; page-0034.tex
\begin{theorem}[线性与二次型求导]\label{reg:entry:33}
设 $A\in\R^{K\times K}$、$x,w\in\R^K$，$A,w$ 不依赖 $x$。按上述排列约定，
\begin{align*}
\frac{\partial(w'x)}{\partial x}&=\frac{\partial(x'w)}{\partial x}=w,\\
\frac{\partial(x'A)}{\partial x}&=A,\qquad \frac{\partial(Ax)}{\partial x'}=A,\\
\frac{\partial(x'Ax)}{\partial x}&=(A+A')x,\\
\frac{\partial^2(x'Ax)}{\partial x\partial x'}&=A+A'.
\end{align*}
若 $A$ 对称，后两式化为 $2Ax$ 和 $2A$。
\end{theorem}


\chapter{古典线性回归模型}

\begin{introduction}[本章的作用]
\item OLS 用样本矩估计总体线性投影系数。
\item 投影几何解释拟合值、残差与自由度。
\item 外生性和球形方差支持无偏性与 BLUE；条件正态性支持精确推断。
\end{introduction}
总体单个观测的解释向量写为 $X_i\in\R^K$；样本设计矩阵写为 $\XX\in\R^{n\times K}$，第 $i$ 行为 $X_i'$。本章的分布与独立性均明确以 $\XX$ 为条件。

\section{经典模型的假设}

% SOURCE: 假设 3.1.1; page-0036.tex
\begin{proposition}[线性（假设 3.1.1）]\label{reg:entry:34}
设 $(Y_i,X_i')_{i=1}^n$ 为可观测随机样本，
\[Y_i=X_i'\beta+u_i,\quad i=1,\ldots,n,\qquad \beta\in\R^K.\]
$u_i$ 是不可观测扰动，$Y=(Y_1,\ldots,Y_n)'$、$u=(u_1,\ldots,u_n)'$，于是
\[Y=\XX\beta+u,\quad \XX=(X_1,\ldots,X_n)',\quad \XX'\XX=\sum_iX_iX_i'.\]
含截距时 $X_i=(1,X_{i1},\ldots,X_{ik})'$，$K=k+1$。
\end{proposition}

% SOURCE: 假设 3.1.2; page-0037.tex
\begin{proposition}[严格外生性（假设 3.1.2）]\label{reg:entry:35}
对于上述样本模型，假设
\[\E[u_i\mid\XX]=\E[u_i\mid X_1,\ldots,X_n]=0,\quad i=1,\ldots,n.\]
即 $\E[u\mid\XX]=0$。从而 $\E u_i=0$，且在相应乘积可积时 $\E[X_i u_j]=0$、$\Cov(X_i,u_j)=0$ 对所有 $i,j$ 成立。
严格外生性一般强于逐观测 $\E[u_i\mid X_i]=0$；当观测对 $(X_i,u_i)$ 独立同分布时，两条件等价。
\end{proposition}

% SOURCE: 假设 3.1.3; page-0038.tex
\begin{proposition}[秩条件（假设 3.1.3）]\label{reg:entry:36}
假设 $\rank(\XX)=K<n$（随机设计时几乎必然）。则列线性无关，$\rank(\XX'\XX)=K$，$\XX'\XX$ 非奇异、可逆且正定。这排除严格多重共线性，确保 OLS 参数解唯一并留下正残差自由度。
\end{proposition}

% SOURCE: 假设 3.1.4; page-0038.tex
\begin{proposition}[球形误差方差（假设 3.1.4）]\label{reg:entry:37}
假设 $\sigma^2>0$ 且
\[\E[u_i^2\mid\XX]=\sigma^2,\qquad \E[u_i u_j\mid\XX]=0\quad(i\ne j).\]
结合严格外生性，等价于
\[\Var(u\mid\XX)=\E[uu'\mid\XX]=\sigma^2 I_n.\]
这同时要求条件同方差与跨观测条件不相关；不相关本身不等于独立。
\end{proposition}

\section{普通最小二乘估计}

% SOURCE: 定义 3.2.1; page-0040.tex
\begin{definition}[OLS]\label{reg:entry:38}
线性模型的普通最小二乘估计定义为
\[\bb=\argmin_{b\in\R^K}\operatorname{SSE}(b),\qquad \operatorname{SSE}(b)=\sum_i(Y_i-X_i'b)^2=(Y-\XX b)'(Y-\XX b).\]
它用样本平方损失替代总体平方损失。
\end{definition}

% SOURCE: 定理 3.2.1; page-0040.tex
\begin{theorem}[OLS 的显式解]\label{reg:entry:39}
在线性模型与秩条件下，OLS 解存在且唯一，
\[\bb=(\XX'\XX)^{-1}\XX'Y
=\left(\frac1n\sum_iX_iX_i'\right)^{-1}\frac1n\sum_iX_iY_i.\]
正规方程为 $\XX'\XX\bb=\XX'Y$，目标函数的 Hessian 为 $2\XX'\XX\succ0$。
\end{theorem}

% SOURCE: 原文未编号概念或必要公式; page-0041.tex
\begin{definition}[拟合值、残差与方差估计]\label{reg:entry:40}
定义 $\widehat Y=\XX\bb$ 和 $\uh=Y-\widehat Y$。则
\[\uh=u+\XX(\beta-\bb).\]
扰动 $u$ 是模型中不可观测的真实误差，残差 $\uh$ 由估计参数计算，两者不可混同。定义
\[\widehat\sigma^2=\frac{\uh'\uh}{n},\qquad s^2=\frac{\uh'\uh}{n-K}.\]
后者在经典假设下修正自由度，并非任意模型下都无偏。
\end{definition}

% SOURCE: 定理 3.2.2; page-0041.tex
\begin{theorem}[残差的样本正交性]\label{reg:entry:41}
在线性模型与秩条件下，
\[\XX'\uh=0,\qquad \widehat Y'\uh=0.\]
若解释变量含常数项，则 $\sum_i\uh_i=0$，进而 $\overline{\widehat Y}=\bar Y$。这是 OLS 的代数性质，不要求正态性、外生性或同方差。
\end{theorem}

\section{投影几何与拟合优度}

% SOURCE: 原文未编号概念或必要公式; page-0042.tex 至 page-0044.tex
\begin{definition}[投影矩阵与消灭矩阵]\label{reg:entry:42}
若 $\XX$ 满列秩，定义
\[P=\XX(\XX'\XX)^{-1}\XX',\qquad M=I_n-P.\]
$P$ 是到 $\mathcal C(\XX)$ 的正交投影，$M$ 是到该列空间正交补的投影。于是
\[\widehat Y=PY,\qquad \uh=MY=Mu,\qquad P\XX=\XX,\quad M\XX=0,\quad PM=MP=0.\]
\end{definition}

% SOURCE: 定理 3.3.1; page-0042.tex
\begin{theorem}[投影矩阵的性质]\label{reg:entry:43}
对任意满列秩 $\XX\in\R^{n\times K}$，$n\geq K$，投影矩阵 $P$ 满足
\[P'=P,\quad P'P=P^2=P,\quad \tr(P)=K,\quad\rank(P)=K.\]
$P$ 的特征值为 $K$ 个 $1$ 和 $n-K$ 个 $0$。
\end{theorem}

% SOURCE: 定理 3.3.2; page-0043.tex
\begin{theorem}[消灭矩阵的性质]\label{reg:entry:44}
对任意满列秩 $\XX\in\R^{n\times K}$，$n\geq K$，$M=I_n-P$ 满足
\[M'=M,\quad M'M=M^2=M,\quad \tr(M)=n-K,\quad\rank(M)=n-K.\]
同时 $M\XX=0$。
\end{theorem}

% SOURCE: 原文未编号概念或必要公式; page-0044.tex 至 page-0045.tex
\begin{proposition}[正交分解与平方和]\label{reg:entry:45}
满列秩时 $Y=\widehat Y+\uh$ 且 $\widehat Y'\uh=0$，所以 $Y'Y=\widehat Y'\widehat Y+\uh'\uh$。
若含截距，中心化后有
\begin{align*}
\operatorname{TSS}&=\sum_i(Y_i-\bar Y)^2,\quad
\operatorname{ESS}=\sum_i(\widehat Y_i-\bar Y)^2,\\
\operatorname{RSS}&=\uh'\uh,\qquad \operatorname{TSS}=\operatorname{ESS}+\operatorname{RSS}.
\end{align*}
对应子空间维数（自由度）为 $n-1=(K-1)+(n-K)$。含截距条件不能省略。
\end{proposition}

% SOURCE: 定义 3.3.1; page-0045.tex
\begin{definition}[决定系数]\label{reg:entry:46}
含截距且 $\operatorname{TSS}>0$ 时，决定系数定义为
\[\mathcal R^2=\frac{\operatorname{ESS}}{\operatorname{TSS}}=1-\frac{\operatorname{RSS}}{\operatorname{TSS}},\qquad0\leq\mathcal R^2\leq1.\]
它度量样本中被线性拟合解释的中心化平方和比例。无截距时这两个表达式一般不等价。
\end{definition}

% SOURCE: 定理 3.3.3; page-0045.tex
\begin{theorem}[加入变量后的决定系数]\label{reg:entry:47}
在同一响应向量、同一样本、含截距且 $\operatorname{TSS}>0$ 的两个嵌套 OLS 模型中，较大模型在原 $k$ 个解释变量外增加 $q\geq1$ 个变量；两设计矩阵满足相应秩条件。则
\[\operatorname{RSS}_{2}\leq\operatorname{RSS}_{1},\qquad \mathcal R_2^2\geq\mathcal R_1^2.\]
这只说明训练样本拟合不会变差，不保证预测性能提高。
\end{theorem}

% SOURCE: 原文未编号概念或必要公式; page-0046.tex
\begin{definition}[调整决定系数]\label{reg:entry:48}
含截距、$n>K$ 且 $\operatorname{TSS}>0$ 时，定义
\[\overline{\mathcal R}^2=1-\frac{\operatorname{RSS}/(n-K)}{\operatorname{TSS}/(n-1)}
=1-\frac{n-1}{n-K}(1-\mathcal R^2).\]
它通过自由度惩罚解释变量数量，增加变量不保证其上升；可为负。
\end{definition}

\section{OLS 的统计性质}

% SOURCE: 定理 3.4.1; page-0047.tex
\begin{theorem}[无偏性]\label{reg:entry:49}
在线性模型、严格外生性和秩条件下，
\[\E[\bb\mid\XX]=\beta.\]
若无条件期望存在，则 $\E\bb=\beta$。所有无条件矩结论均需相应可积性；随机矩阵逆的可积性并非仅由逐样本满秩保证。
\end{theorem}

% SOURCE: 定理 3.4.2; page-0047.tex
\begin{theorem}[OLS 条件协方差]\label{reg:entry:50}
在线性模型、严格外生性、秩条件和球形误差方差条件下，
\[V_{\bb}:=\Var(\bb\mid\XX)=\sigma^2(\XX'\XX)^{-1}.\]
\end{theorem}

% SOURCE: 定义 3.4.1; page-0047.tex
\begin{definition}[线性无偏估计量]\label{reg:entry:51}
若估计量 $\widetilde\beta=L Y$，其中 $L\in\R^{K\times n}$ 可依赖于 $\XX$，但给定 $\XX$ 后不依赖于 $Y$，称其对 $Y$ 为线性估计。
若对所有 $\beta$，$\E[\widetilde\beta\mid\XX]=\beta$，称线性无偏估计。在外生模型下，这等价于 $L\XX=I_K$。
\end{definition}

% SOURCE: 定理 3.4.3; page-0048.tex
\begin{theorem}[Gauss--Markov]\label{reg:entry:52}
在线性模型、严格外生性、秩条件及球形误差方差条件下，OLS 是最优线性无偏估计（BLUE）。对于任意线性无偏估计量 $\widetilde\beta$，
\[\Var(\widetilde\beta\mid\XX)-\Var(\bb\mid\XX)\succeq0.\]
“最小”指协方差矩阵的半正定顺序；不声称在全部非线性或有偏估计量中最优，也不要求正态误差。
\end{theorem}

% SOURCE: 定理 3.4.4; page-0049.tex
\begin{theorem}[残差的矩性质]\label{reg:entry:53}
在线性模型、严格外生性、秩条件及球形误差方差条件下，
\[\E[\uh\mid\XX]=0,\quad \Var(\uh\mid\XX)=\sigma^2M,\quad \Cov(\bb,\uh\mid\XX)=0.\]
条件不相关尚不能推出条件独立。
\end{theorem}

% SOURCE: 定理 3.4.5; page-0050.tex
\begin{theorem}[方差估计的无偏性]\label{reg:entry:54}
在线性模型、严格外生性、秩条件及球形误差方差条件下，
\[\E[\uh'\uh\mid\XX]=(n-K)\sigma^2,\quad \E[s^2\mid\XX]=\sigma^2.\]
所以 $\E[\widehat\sigma^2\mid\XX]=(n-K)\sigma^2/n$；OLS 协方差的无偏估计为
\[\widehat V_{\bb}=s^2(\XX'\XX)^{-1},\qquad \E[\widehat V_{\bb}\mid\XX]=V_{\bb}.\]
\end{theorem}

\section{正态模型的精确分布}

% SOURCE: 假设 3.5.1; page-0050.tex
\begin{proposition}[条件正态性（假设 3.5.1）]\label{reg:entry:55}
假设 $u\mid\XX\sim N(0,\sigma^2I_n)$，$\sigma^2>0$。该条件推出严格外生性和球形方差，且给定 $\XX$ 时各扰动独立。
其条件密度和似然为
\[f(u\mid\XX)=(2\pi\sigma^2)^{-n/2}\exp[-u'u/(2\sigma^2)],\]
\[\ell(b,\sigma^2)=-\frac n2\log(2\pi\sigma^2)-\frac{(Y-\XX b)'(Y-\XX b)}{2\sigma^2}.\]
在满列秩且 RSS 正的情形，极大似然解为 $\widehat b_{\rm MLE}=\bb$、$\widehat\sigma^2_{\rm MLE}=\operatorname{RSS}/n$。
\end{proposition}

% SOURCE: 定理 3.5.1; page-0051.tex
\begin{theorem}[OLS 的条件正态分布]\label{reg:entry:56}
在线性模型、秩条件和条件正态性下，$n>K$，
\[\bb\mid\XX\sim N\!\left(\beta,\sigma^2(\XX'\XX)^{-1}\right).\]
\end{theorem}

% SOURCE: 定理 3.5.2; page-0052.tex
\begin{theorem}[残差分布及条件独立性]\label{reg:entry:57}
在线性模型、秩条件和条件正态性下，$n>K$，
\[\uh\mid\XX\sim N(0,\sigma^2M).\]
给定 $\XX$ 时，$\uh$ 与 $\bb$ 独立。因为 $M$ 秩为 $n-K<n$，这个残差正态分布是退化的。
\end{theorem}

% SOURCE: 定理 3.5.3; page-0052.tex
\begin{theorem}[误差方差估计的分布]\label{reg:entry:58}
在线性模型、秩条件和条件正态性下，$n>K$，
\[\frac{(n-K)s^2}{\sigma^2}\biggm|\XX\sim\chi^2_{n-K}.\]
给定 $\XX$ 时，$s^2$ 与 $\bb$ 独立。这里的条件独立不能一概改成无条件独立。
\end{theorem}

\section{线性限制的假设检验}

% SOURCE: 原文未编号概念或必要公式; page-0053.tex 至 page-0054.tex
\begin{definition}[线性限制与拒绝域]\label{reg:entry:59}
令 $R\in\R^{J\times K}$ 满行秩，$r\in\R^J$，$1\leq J\leq K$，检验
\[H_0:R\beta=r,\qquad H_1:R\beta\ne r.\]
检验用原假设下的统计量分布确定显著性水平 $\alpha$ 的拒绝域。落在拒绝域时拒绝 $H_0$；没有拒绝不能当作证明 $H_0$ 为真。令 $Q=(\XX'\XX)^{-1}$，则 $RQR'\succ0$。
\end{definition}

% SOURCE: 定理 3.6.1; page-0054.tex
\begin{theorem}[限制估计量的正态分布]\label{reg:entry:60}
在线性模型、秩条件和条件正态性下，若 $H_0:R\beta=r$ 成立，且 $R$ 满行秩，
\[R\bb-r\mid\XX\sim N(0,\sigma^2RQR').\]
\end{theorem}

% SOURCE: 定理 3.6.2; page-0054.tex
\begin{theorem}[单个限制的 t 检验]\label{reg:entry:61}
在线性模型、秩条件、条件正态性和 $H_0:R\beta=r$ 下，$n>K$、$J=1$ 且 $R\ne0$，
\[T=\frac{R\bb-r}{\sqrt{s^2RQR'}}\biggm|\XX\sim t_{n-K}.\]
双侧水平 $\alpha$ 检验在 $|T|>t_{n-K,1-\alpha/2}$ 时拒绝。对单个系数的置信区间为
\[\bb_j\ \pm\ t_{n-K,1-\alpha/2}\sqrt{s^2Q_{jj}}.\]
\end{theorem}

% SOURCE: 原文未编号概念或必要公式; page-0055.tex
\begin{proposition}[联合限制的 F 统计量]\label{reg:entry:62}
在相同正态模型条件下，$R$ 满行秩，且 $H_0$ 成立，则
\[F=\frac{(R\bb-r)'(RQR')^{-1}(R\bb-r)}{J s^2}\biggm|\XX\sim F_{J,n-K}.\]
水平 $\alpha$ 检验在 $F>F_{J,n-K,1-\alpha}$ 时拒绝。$J=1$ 时 $F=T^2$。
\end{proposition}

% SOURCE: 定理 3.6.3; page-0056.tex
\begin{theorem}[约束回归形式的 F 检验]\label{reg:entry:63}
在线性模型、秩条件、条件正态性及 $H_0:R\beta=r$ 下，$n>K$，$R$ 有 $J$ 个独立限制。令 $\widetilde u$ 为约束最小二乘残差，$\uh$ 为无约束 OLS 残差，则
\[F_1=\frac{(\widetilde u'\widetilde u-\uh'\uh)/J}{\uh'\uh/(n-K)}\biggm|\XX\sim F_{J,n-K}.\]
它与上面的二次型形式相等。
\end{theorem}

% SOURCE: 原文未编号概念或必要公式; page-0057.tex
\begin{proposition}[所有斜率的联合检验]\label{reg:entry:64}
含截距模型有 $K-1\geq1$ 个斜率。检验所有斜率为零时，约束拟合为 $\bar Y$，故约束 RSS 等于 TSS。若 TSS 正，则
\[F=\frac{\mathcal R^2/(K-1)}{(1-\mathcal R^2)/(n-K)}.\]
在上述正态模型与原假设下，$F\mid\XX\sim F_{K-1,n-K}$。
\end{proposition}

\section{拾遗：随机矩阵的矩}

% SOURCE: 原文未编号概念或必要公式; page-0058.tex
\begin{definition}[随机向量协方差]\label{reg:entry:65}
随机矩阵期望逐元素定义，要求每个元素可积。对于二阶矩有限的向量 $X,Y$，
\[\Cov(X,Y)=\E[(X-\E X)(Y-\E Y)'],\quad \Var(X)=\Cov(X,X).\]
固定矩阵 $A,B$ 满足维度相容时，$\Cov(AX,BY)=A\Cov(X,Y)B'$。
\end{definition}

% SOURCE: 定理 3.7.1; page-0058.tex
\begin{theorem}[期望的线性性质]\label{reg:entry:66}
若随机矩阵或向量 $X,Y$ 的元素可积，$A,B$ 为维度相容的常数矩阵，则
\[\E[AX+BY]=A\E[X]+B\E[Y].\]
对于可积方阵，$\tr(\E X)=\E(\tr X)$。
\end{theorem}

% SOURCE: 定理 3.7.2; page-0058.tex
\begin{theorem}[协方差矩阵的性质]\label{reg:entry:67}
若 $X\in\R^K$ 且 $\E[X'X]<\infty$，记 $\mu=\E X$，$A$ 为常数矩阵、$b$ 为常数向量，则
\[\Var(X)\succeq0,\quad \Var(X)=\E[XX']-\mu\mu',\quad \Var(AX+b)=A\Var(X)A'.\]
\end{theorem}

% SOURCE: 原文未编号概念或必要公式; 补充工具公式：由协方差定义与迹性质得到
\begin{proposition}[二次型期望的必要公式]\label{reg:entry:68}
设 $X$ 二阶矩有限，$\mu=\E X$、$\Sigma=\Var(X)$，$A$ 为常数方阵，则
\[\E[X'AX]=\tr(A\Sigma)+\mu'A\mu.\]
此式把 RSS 的期望转化为投影矩阵的迹，是自由度修正的工具。
\end{proposition}

\section{拾遗：多元正态与二次型}

% SOURCE: 定义 3.8.1; page-0059.tex
\begin{definition}[多元正态分布]\label{reg:entry:69}
若 $Z\sim N(0,I_m)$，$Y=\mu+BZ\in\R^d$，则称 $Y$ 多元正态，记 $Y\sim N(\mu,\Sigma)$，其中 $\Sigma=BB'$，$\E Y=\mu$、$\Var(Y)=\Sigma$。
若 $\Sigma\succ0$，其密度为
\[f(y)=(2\pi)^{-d/2}|\Sigma|^{-1/2}\exp[-(y-\mu)'\Sigma^{-1}(y-\mu)/2].\]
若 $\Sigma$ 奇异，分布仍可由仿射表示定义，但不能使用这条相对于 $\R^d$ Lebesgue 测度的密度公式。
\end{definition}

% SOURCE: 定理 3.8.1; page-0059.tex
\begin{theorem}[正态向量的仿射变换]\label{reg:entry:70}
若 $Y\sim N(\mu,\Sigma)$，$a,C$ 为维度相容的常数向量和矩阵，则
\[a+CY\sim N(a+C\mu,C\Sigma C').\]
\end{theorem}

% SOURCE: 定理 3.8.2; page-0060.tex
\begin{theorem}[分块正态分布]\label{reg:entry:71}
设 $Y=(Y_1',Y_2')'$ 联合正态，均值分块为 $(\mu_1',\mu_2')'$，协方差分块为 $(\Sigma_{ij})_{i,j=1,2}$。则
\begin{enumerate}
\item $Y_1\sim N(\mu_1,\Sigma_{11})$。
\item 若 $\Sigma_{22}$ 可逆，则
\[Y_1\mid Y_2\sim N\!\left(\mu_1+\Sigma_{12}\Sigma_{22}^{-1}(Y_2-\mu_2),\,
\Sigma_{11}-\Sigma_{12}\Sigma_{22}^{-1}\Sigma_{21}\right).\]
\item 若 $\Sigma_{12}=0$，则 $Y_1,Y_2$ 独立；联合正态下此条件也必要。
\end{enumerate}
\end{theorem}

% SOURCE: 定理 3.8.3; page-0055.tex
\begin{theorem}[正态平方和与卡方分布]\label{reg:entry:72}
\begin{enumerate}
\item 若 $Y\sim N(0,I_n)$，则 $Y'Y\sim\chi_n^2$。
\item 若 $Y\sim N(0,\Sigma)$ 且 $\Sigma\succ0$，则 $Y'\Sigma^{-1}Y\sim\chi_n^2$。
\item 若 $Y\sim N(\mu,I_n)$，则 $Y'Y\sim\chi_n^2(\lambda)$，$\lambda=\mu'\mu$。
\item 若 $Y\sim N(\mu,\Sigma)$ 且 $\Sigma\succ0$，则 $Y'\Sigma^{-1}Y\sim\chi_n^2(\lambda)$，$\lambda=\mu'\Sigma^{-1}\mu$。
\end{enumerate}
非中心参数统一采用平方范数惯例，不带 $1/2$ 因子。
\end{theorem}

% SOURCE: 定理 3.8.4; page-0060.tex
\begin{theorem}[正态投影二次型]\label{reg:entry:73}
设 $Q$ 是非随机 $n\times n$ 对称幂等矩阵，$\rank Q=r\leq n$。若 $Y\sim N(0,I_n)$，则
\[Y'QY\sim\chi_r^2.\]
更一般地，若 $Z\sim N(\mu,I_n)$，则 $Z'QZ\sim\chi_r^2(\lambda)$，$\lambda=\mu'Q\mu$。$r=0$ 时按恒为零的退化分布理解。
\end{theorem}

% SOURCE: 定理 3.8.5; page-0061.tex
\begin{theorem}[线性形式与二次型的独立性]\label{reg:entry:74}
若 $Y\sim N(0,I_n)$，$A\in\R^{m\times n}$，$B\in\R^{n\times n}$ 为对称常数矩阵，且 $AB=0$，则 $AY$ 与 $Y'BY$ 独立。
\end{theorem}

% SOURCE: 定理 3.8.6; page-0061.tex
\begin{theorem}[两个投影二次型的独立性]\label{reg:entry:75}
若 $Y\sim N(0,I_n)$，$A,B$ 为 $n\times n$ 对称幂等常数矩阵，则
\[Y'AY\ \text{与}\ Y'BY\ \text{独立}\quad\Longleftrightarrow\quad AB=0.\]
\end{theorem}

\section{拾遗：约束最小二乘回归}

% SOURCE: 原文未编号概念或必要公式; page-0061.tex 至 page-0062.tex
\begin{definition}[CLS 及其显式解]\label{reg:entry:76}
令 $Q=(\XX'\XX)^{-1}$，$R\in\R^{J\times K}$ 满行秩，$1\leq J\leq K$，$r\in\R^J$。在满列秩条件下定义
\[\widetilde\beta=\argmin_{Rb=r}(Y-\XX b)'(Y-\XX b).\]
解及相关矩阵为
\begin{align*}
\widetilde\beta&=\bb-QR'(RQR')^{-1}(R\bb-r),\\
A&=QR'(RQR')^{-1}RQ,\qquad D=(Q-A)\XX'.
\end{align*}
$RQR'\succ0$；只有真实参数满足 $R\beta=r$ 时，才有 $\widetilde\beta=\beta+Du$。
\end{definition}

% SOURCE: 定理 3.9.1; page-0062.tex
\begin{theorem}[CLS 的无偏性]\label{reg:entry:77}
在线性模型、严格外生性、秩条件及真实限制 $R\beta=r$ 下，$R$ 满行秩，则
\[\E[\widetilde\beta\mid\XX]=\beta.\]
错误限制通常导致有偏，不能省去真实限制条件。
\end{theorem}

% SOURCE: 定理 3.9.2; page-0063.tex
\begin{theorem}[CLS 的条件协方差]\label{reg:entry:78}
在线性模型、严格外生性、秩条件、球形误差方差及真实限制 $R\beta=r$ 下，$R$ 满行秩，则
\[V_{\widetilde\beta}:=\Var(\widetilde\beta\mid\XX)=\sigma^2DD'=\sigma^2(Q-A).\]
\end{theorem}

% SOURCE: 定理 3.9.3; page-0063.tex
\begin{theorem}[真实限制带来的方差改进]\label{reg:entry:79}
在线性模型、严格外生性、秩条件、球形误差方差及真实限制 $R\beta=r$ 下，$R$ 满行秩，则
\[V_{\bb}-V_{\widetilde\beta}=\sigma^2 A\succeq0.\]
限制利用额外的真实信息降低方差，与 Gauss--Markov 在无约束参数空间上的最优性并不冲突。
\end{theorem}

% SOURCE: 原文未编号概念或必要公式; page-0063.tex 至 page-0064.tex
\begin{definition}[CLS 残差与方差估计]\label{reg:entry:80}
令 $\widetilde u=Y-\XX\widetilde\beta$。在真实限制 $R\beta=r$ 下，
\[\widetilde u=Cu,\qquad C=M+\XX A\XX'.\]
$C$ 对称幂等，$\tr C=\rank C=n-K+J$。定义
\[s_{\rm cls}^2=\frac{\widetilde u'\widetilde u}{n-K+J},\qquad \widehat V_{\widetilde\beta}=s_{\rm cls}^2(Q-A).\]
\end{definition}

% SOURCE: 定理 3.9.4; page-0064.tex
\begin{theorem}[CLS 方差估计的无偏性]\label{reg:entry:81}
在线性模型、严格外生性、秩条件、球形误差方差及真实限制 $R\beta=r$ 下，$R$ 满行秩，则
\[\E[s_{\rm cls}^2\mid\XX]=\sigma^2,\qquad \E[\widehat V_{\widetilde\beta}\mid\XX]=V_{\widetilde\beta}.\]
\end{theorem}

% SOURCE: 定理 3.9.5; page-0065.tex
\begin{theorem}[CLS 的条件抽样分布]\label{reg:entry:82}
在线性模型、秩条件、条件正态性及真实限制 $R\beta=r$ 下，$R$ 满行秩，则
\[\widetilde\beta\mid\XX\sim N\!\left(\beta,\sigma^2(Q-A)\right),\qquad
\frac{(n-K+J)s_{\rm cls}^2}{\sigma^2}\biggm|\XX\sim\chi^2_{n-K+J}.\]
参数分布一般退化，反映 $J$ 个确定限制。
\end{theorem}


\begingroup
\chapter*{第二章作业：条件期望函数}
\addcontentsline{toc}{chapter}{第二章作业：条件期望函数}
\markboth{第二章作业：条件期望函数}{第二章作业：条件期望函数}
\setcounter{section}{0}
\renewcommand{\thesection}{\arabic{section}}
\renewcommand{\thechapter}{2}
\renewcommand{\theHsection}{reg.homework.2.\arabic{section}}

\section{第一次作业}\label{hw:reg:assignment:1}

阅读至第 2.1.2 节前；整理相关系数、重期望公式与最优预测的证明，完成练习 1--6。相关系数需正方差；最优平方损失预测需 $\E Y^2<\infty$。

\subsection*{课程笔记}

\begin{proposition}[阅读笔记：协方差与相关系数]\label{hw:reg:note:1:1}
若 $X,Y$ 二阶矩有限，协方差定义为
\[
\Cov(X,Y)=\E[(X-\E X)(Y-\E Y)]
=\E(XY)-\E X\,\E Y.
\]
它刻画两个变量的线性共同变动，具有
\[
\Cov(aX+b,cY+d)=ac\,\Cov(X,Y),\qquad
|\Cov(X,Y)|\leq\sqrt{\Var(X)\Var(Y)}.
\]
当 $\Var(X),\Var(Y)>0$ 时，相关系数为
\[
\rho_{XY}=\frac{\Cov(X,Y)}
{\sqrt{\Var(X)\Var(Y)}}\in[-1,1].
\]
$|\rho|=1$ 当且仅当 $Y-\E Y=c(X-\E X)$ 几乎处处成立。独立必然推出零相关，但零相关通常不能推出独立；若 $(X,Y)$ 联合正态，则零相关与独立等价。

样本相关系数是
\[
R=\frac{\sum_{i=1}^n(X_i-\bar X)(Y_i-\bar Y)}
{\sqrt{
\left[\sum_{i=1}^n(X_i-\bar X)^2\right]
\left[\sum_{i=1}^n(Y_i-\bar Y)^2\right]
}},
\]
它对平移与正比例缩放不变。相关性只表示线性关系强弱，不自动表示因果关系。
\end{proposition}

\begin{proposition}[阅读笔记：重期望公式及证明]\label{hw:reg:note:1:2}
若 $\E|Y|<\infty$，则
\[
\E\{\E(Y\mid X)\}=\E Y.
\]
连续情形中，设联合密度为 $f(x,y)$，边缘密度为 $f_X(x)$，条件密度为 $f(y\mid x)$，则
\begin{align*}
\E\{\E(Y\mid X)\}
&=\int\left[\int y f(y\mid x)\,dy\right]f_X(x)\,dx\\
&=\iint y f(x,y)\,dy\,dx\\
&=\int y f_Y(y)\,dy=\E Y.
\end{align*}
更一般地，只要 $\E|G(X,Y)|<\infty$，
\[
\E\{\E[G(X,Y)\mid X]\}=\E G(X,Y).
\]
证明的本质是：先在每个给定 $X=x$ 的子总体内平均，再按 $X$ 的分布加权，恰好恢复总体平均。
\end{proposition}

\begin{proposition}[阅读笔记：最优预测及证明]\label{hw:reg:note:1:3}
用平方损失预测 $Y$，对任意可测且平方可积的函数 $g(X)$，定义
\[
\operatorname{MSE}(g)=\E[(Y-g(X))^2].
\]
令 $\mu(X)=\E(Y\mid X)$。把误差分解为
\[
Y-g(X)=\{Y-\mu(X)\}+\{\mu(X)-g(X)\}.
\]
平方并取期望。交叉项等于
\[
\E\!\left[\{\mu(X)-g(X)\}\E\{Y-\mu(X)\mid X\}\right]=0.
\]
因此
\[
\operatorname{MSE}(g)
=\E[(Y-\mu(X))^2]+\E[(\mu(X)-g(X))^2].
\]
第一项与 $g$ 无关，第二项非负，所以 $g(X)=\mu(X)$ 时均方误差最小。故条件期望是平方损失下的最优预测。
\end{proposition}


\Needspace{10\baselineskip}
\subsection*{练习 1：二元正态与独立性}
\begin{homework}[练习 1]
\textsf{来源：《回归分析 60 题》已有原题第 2 页文件。}\par
设 $(X,Y)$ 服从二元正态分布 $N(\mu_1,\mu_2,\sigma_1^2,\sigma_2^2,\rho)$，其中 $\sigma_1,\sigma_2>0$、$|\rho|<1$。证明 $X,Y$ 的相关系数为 $\rho_{XY}=\rho$；在二元正态条件下，$\rho=0$ 是 $X,Y$ 相互独立的充要条件。
\end{homework}
\begin{solution}
二元正态的协方差矩阵为
\[
\Sigma=\begin{pmatrix}
\sigma_1^2&\rho\sigma_1\sigma_2\\
\rho\sigma_1\sigma_2&\sigma_2^2
\end{pmatrix}.
\]
所以 $\Cov(X,Y)=\rho\sigma_1\sigma_2$，从而 $\rho_{XY}=\rho$。独立必然推出协方差为零；反过来，在联合正态条件下，$\rho=0$ 使联合密度分解为两个边缘密度之积，因此 $X,Y$ 独立。
\end{solution}
\begin{dxtips}[二元正态与独立性]\label{hw:reg:tip:1}
零协方差推出独立必须使用联合正态条件；一般随机变量只有独立推出零协方差。
\end{dxtips}

\Needspace{10\baselineskip}
\subsection*{练习 2：Cauchy--Schwarz 及等号条件}
\begin{homework}[练习 2]
\textsf{来源：《回归分析 60 题》已有原题第 2 页文件。}\par
设 $X,Y$ 二阶矩有限。证明柯西--施瓦茨不等式，讨论等号成立的条件。需区分 $|\E(XY)|\leq\sqrt{\E X^2\E Y^2}$ 和 $\E|XY|\leq\sqrt{\E X^2\E Y^2}$ 两种形式。
\end{homework}
\begin{solution}
若 $\E(Y^2)>0$，对任意 $t$，
\[
0\leq \E(X-tY)^2
=\E X^2-2t\E(XY)+t^2\E Y^2.
\]
取 $t=\E(XY)/\E(Y^2)$ 即得不等式。当 $\E Y^2>0$ 时，等号成立当且仅当 $X=tY$ 几乎处处。若 $\E Y^2=0$，则 $Y=0$ 几乎处处，不等式恒取等号，无须 $X=0$。统一条件是存在不全为零的常数 $a,b$ 使 $aX+bY=0$ 几乎处处。
对绝对乘积形式，把同一证明用于 $|X|,|Y|$，得到 $\E|XY|\leq\sqrt{\E X^2\E Y^2}$；非退化时其等号条件为 $|X|=c|Y|$ 几乎处处。
\end{solution}
\begin{dxtips}[Cauchy--Schwarz 及等号条件]\label{hw:reg:tip:2}
二阶矩为零的变量几乎处处为零。非退化时对 $\E(X-tY)^2$ 最小化；绝对乘积形式则把 $X,Y$ 换成 $|X|,|Y|$。
\end{dxtips}

\Needspace{10\baselineskip}
\subsection*{练习 3：期望的绝对值}
\begin{homework}[练习 3]
\textsf{来源：《回归分析 60 题》已有原题第 2 页文件。}\par
证明期望不等式：对可积随机变量 $X$，有 $|\E X|\leq\E|X|$。原题提示：利用詹森不等式。
\end{homework}
\begin{solution}
$\phi(x)=|x|$ 为凸函数，由 Jensen 不等式
\[
|\E X|=\phi(\E X)\leq\E\phi(X)=\E|X|.
\]
\end{solution}
\begin{dxtips}[期望的绝对值]\label{hw:reg:tip:3}
先检查可积性；绝对值是凸函数，直接应用 Jensen。
\end{dxtips}

\Needspace{10\baselineskip}
\subsection*{练习 4：线性条件均值下的混合矩}
\begin{homework}[练习 4]
\textsf{来源：《回归分析 60 题》已有原题第 2 页文件。}\par
若 $\E[Y\mid X]=a+bX$，请证明 $\E[YX]$ 是 $X$ 的矩的函数。假定相关期望存在，例如 $X,Y$ 二阶矩有限。
\end{homework}
\begin{solution}
\[
\E(YX)=\E\{X\E(Y\mid X)\}
=a\E X+b\E(X^2).
\]
\end{solution}
\begin{dxtips}[线性条件均值下的混合矩]\label{hw:reg:tip:4}
先对 X 条件化，再把可测的 X 提出。
\end{dxtips}

\Needspace{10\baselineskip}
\subsection*{练习 5：条件混合矩与重期望}
\begin{homework}[练习 5]
\textsf{来源：《回归分析 60 题》已有原题第 2 页文件。}\par
如果 $\E[Y\mid X]=a+bX$，请证明 $\E[YX]$ 是 $X$ 的矩的函数。假定相关期望存在，例如 $X,Y$ 二阶矩有限。此题与第 4 题重复，沿用原题编号分别保留。
\end{homework}
\begin{solution}
因为 $X$ 对 $\sigma(X)$ 可测，
\[
\E(YX\mid X)=X\E(Y\mid X)=aX+bX^2.
\]
两边取无条件期望即得 $\E(YX)=a\E X+b\E(X^2)$。
\end{solution}
\begin{dxtips}[条件混合矩与重期望]\label{hw:reg:tip:5}
原题重复不意味着漏题：此处从条件混合矩写起，再用重期望，保留另一种书写顺序。
\end{dxtips}

\Needspace{10\baselineskip}
\subsection*{练习 6：广义重期望}
\begin{homework}[练习 6]
\textsf{来源：《回归分析 60 题》已有原题第 2 页文件。}\par
对任意可测函数 $G(X,Y)$，若 $\E|G(X,Y)|<\infty$，证明
\[\E_X\bigl[\E_Y[G(X,Y)\mid X]\bigr]=\E_{X,Y}[G(X,Y)].\]
\end{homework}
\begin{solution}
根据条件期望的定义，对全集 $\Omega\in\sigma(X)$ 有
\[
\int_\Omega \E(G\mid X)\,dP=\int_\Omega G\,dP.
\]
左、右两边正是 $\E\{\E(G\mid X)\}$ 与 $\E G$，所以结论成立。
\end{solution}
\begin{dxtips}[广义重期望]\label{hw:reg:tip:6}
条件期望的定义直接包含对全集积分的恒等式，证明不限于存在密度的情形。
\end{dxtips}

\subsection*{知识点定位}\begin{center}\begin{tabular}{ll}\toprule 题号 & 核心方法\\\midrule

1 & 二元正态与独立性 \\

2 & Cauchy--Schwarz 及等号条件 \\

3 & 期望的绝对值 \\

4 & 线性条件均值下的混合矩 \\

5 & 条件混合矩与重期望 \\

6 & 广义重期望 \\

\bottomrule\end{tabular}\end{center}
\section{第二次作业}\label{hw:reg:assignment:2}

阅读至第 2.3 节前；整理条件方差、CEF 误差和最优线性预测，完成练习 8、9、10、11、18。各二阶矩公式按相应平方可积条件使用。

\subsection*{课程笔记}

\begin{proposition}[阅读笔记：条件方差与全方差公式]\label{hw:reg:note:2:1}
条件方差为
\[
\Var(Y\mid X)=\E(Y^2\mid X)-\{\E(Y\mid X)\}^2.
\]
令 $\mu(X)=\E(Y\mid X)$，则
\[
Y-\E Y=\{Y-\mu(X)\}+\{\mu(X)-\E Y\}.
\]
两部分正交，因为 $\E[Y-\mu(X)\mid X]=0$。因此
\[
\Var(Y)=\E\{\Var(Y\mid X)\}
+\Var\{\E(Y\mid X)\}.
\]
第一项是无法由 $X$ 解释的平均波动，第二项是不同 $X$ 子总体条件均值之间的波动。
\end{proposition}

\begin{proposition}[阅读笔记：CEF 误差与正交性]\label{hw:reg:note:2:2}
把
\[
Y=\mu(X)+e,\qquad \mu(X)=\E(Y\mid X)
\]
称为条件期望函数分解。由定义
\[
\E(e\mid X)=0.
\]
于是对任意满足 $\E|h(X)e|<\infty$ 的可测函数 $h(X)$，
\[
\E[h(X)e]=\E\{h(X)\E(e\mid X)\}=0.
\]
因此 CEF 误差不仅与 $X$ 正交，也与 $X$ 的任意可测变换正交。需要注意：$\E(Xe)=0$ 只是一个无条件矩条件，通常不能反推出 $\E(e\mid X)=0$。
\end{proposition}

\begin{proposition}[阅读笔记：线性 CEF 与最优线性预测]\label{hw:reg:note:2:3}
若 $\E(Y\mid X)=X'\beta$，则条件均值被正确设定为线性函数。即使真实 CEF 非线性，也可以在所有线性函数中寻找
\[
\beta=\arg\min_b\E[(Y-X'b)^2].
\]
一阶条件为
\[
\E[X(Y-X'\beta)]=0,
\]
若 $\E(XX')$ 可逆，则
\[
\beta=\{\E(XX')\}^{-1}\E(XY).
\]
该 $X'\beta$ 是最优线性预测。线性投影误差只保证 $\E(Xe)=0$；只有模型正确设定时，才进一步有 $\E(e\mid X)=0$。
\end{proposition}


\Needspace{10\baselineskip}
\subsection*{练习 8：回归误差的高阶矩}
\begin{homework}[练习 8]
\textsf{来源：《回归分析 60 题》已有原题第 2 页文件。}\par
在线性模型 $Y=X'\beta+e$ 下，采用有限二阶矩及 $Q_{XX}=\E[XX']\succ0$ 的总体假设（原习题称“假设 2.1”，讲义对应假设 2.2.1）。证明若
\[\E|Y|^r<\infty,\qquad\E\|X\|^r<\infty,\qquad r\geq2,\]
则 $\E|e|^r<\infty$。原题提示：使用 Minkowski 不等式，对于 $p\geq1$，
\[(\E|U+V|^p)^{1/p}\leq(\E|U|^p)^{1/p}+(\E|V|^p)^{1/p}.\]
\end{homework}
\begin{solution}
由 $e=Y-X'\beta$、Minkowski 不等式和 $|X'\beta|\leq\|X\|\|\beta\|$，
\[
(\E|e|^r)^{1/r}
\leq(\E|Y|^r)^{1/r}
+\|\beta\|(\E\|X\|^r)^{1/r}<\infty.
\]
因此 $\E|e|^r<\infty$。
\end{solution}
\begin{dxtips}[回归误差的高阶矩]\label{hw:reg:tip:8}
先用向量 Cauchy--Schwarz 把 $|X^{\mathsf T}\beta|$ 控制为 $\|X\|\|\beta\|$，再用 Minkowski。该矩界本身不依赖投影最优性。
\end{dxtips}

\Needspace{10\baselineskip}
\subsection*{练习 9：乘法异方差模型}
\begin{homework}[练习 9]
\textsf{来源：《回归分析 60 题》已有原题第 2 页文件。}\par
假设 $Y=\beta_0+\beta_1X_1+|X_1|e$，其中 $\E X_1=0$、$\Var(X_1)=\sigma_{X_1}^2>0$，$\E e=0$、$\Var(e)=\sigma_e^2>0$，$e$ 与 $X_1$ 相互独立，$\beta_0,\beta_1$ 是未知参数。
\begin{enumerate}
\item 求 $\E[Y\mid X_1]$。
\item 求 $\Var(Y\mid X_1)$。
\item 证明 $\beta_1=0$ 当且仅当 $\Cov(X_1,Y)=0$。
\end{enumerate}
\end{homework}
\begin{solution}
\[
\E(Y\mid X_1)=\beta_0+\beta_1X_1,
\qquad
\Var(Y\mid X_1)=X_1^2\sigma_e^2.
\]
又因为 $\E[X_1|X_1|e]=\E[X_1|X_1|]\E e=0$，
\[
\Cov(X_1,Y)=\beta_1\Var(X_1)=\beta_1\sigma_{X_1}^2.
\]
$\sigma_{X_1}^2>0$，所以 $\beta_1=0$ 当且仅当 $\Cov(X_1,Y)=0$。
\end{solution}
\begin{dxtips}[乘法异方差模型]\label{hw:reg:tip:9}
条件化后 $|X_1|$ 是常数，所以方差会带 $X_1^2$。斜率与协方差等价的关键是 $\Var(X_1)>0$。
\end{dxtips}

\Needspace{10\baselineskip}
\subsection*{练习 10：遗漏交互项}
\begin{homework}[练习 10]
\textsf{来源：《回归分析 60 题》已有原题第 2 页文件。}\par
总体模型为 $Y=0.8X_1X_2+e$，其中 $X_1,X_2,e$ 相互独立且均服从 $N(0,1)$。令 $X=(1,X_1,X_2)'$。假设利用只含截距和两个主效应的线性回归模型 $Y=X'\beta+u$ 刻画 $Y$ 与 $X$ 的相关关系，有什么问题？原题沿用 $e$ 表示该拟合误差；这里改记 $u$，以免与真实噪声混淆。
\end{homework}
\begin{solution}
真实条件均值为
\[
\E(Y\mid X_1,X_2)=0.8X_1X_2,
\]
它不是 $(1,X_1,X_2)$ 的线性函数。由于各变量独立标准正态，
\[
\E Y=\E(X_1Y)=\E(X_2Y)=0,
\]
所以该线性投影的三个系数全为零。模型遗漏了关键交互项 $X_1X_2$，会把真实可预测结构误判为“没有关系”。\par 又 $\E[XX']=I_3$，故 $\beta^*=0$。若加入交互项 $X_1X_2$，最优系数为 $(0,0,0,0.8)'$，可恢复真实 CEF。
\end{solution}
\begin{dxtips}[遗漏交互项]\label{hw:reg:tip:10}
交互作用可以存在而所有主效应投影系数都为零。加入 $X_1X_2$ 后对参数仍是线性模型。原设定投影 MSE 为 $1.64$，真实 CEF 的 MSE 为 $1$。
\end{dxtips}

\Needspace{10\baselineskip}
\subsection*{练习 11：离散条件矩}
\begin{homework}[练习 11]
\textsf{来源：《回归分析 60 题》已有原题第 3 页文件。}\par
随机变量 $X,Y$ 只能取 $0$ 或 $1$，联合分布如下（行标为 $Y$，列标为 $X$）：
\[\begin{array}{c|cc}&X=0&X=1\\\hline Y=0&0.1&0.2\\Y=1&0.4&0.3\end{array}\]
求 $X=0$ 或 $X=1$ 时 $\E[Y\mid X]$、$\E[Y^2\mid X]$ 和 $\Var(Y\mid X)$。
\end{homework}
\begin{solution}
$P(X=0)=P(X=1)=0.5$，且 $Y^2=Y$。因此
\[
\begin{array}{c|ccc}
&\E(Y\mid X)&\E(Y^2\mid X)&\Var(Y\mid X)\\\midrule
X=0&0.8&0.8&0.16\\
X=1&0.6&0.6&0.24
\end{array}
\]
\end{solution}
\begin{dxtips}[离散条件矩]\label{hw:reg:tip:11}
必须先确认表格方向，再计算每列的边际概率。$Y$ 为 Bernoulli 型取值，故 $Y^2=Y$、条件方差为 $p(1-p)$。
\end{dxtips}

\Needspace{10\baselineskip}
\subsection*{练习 18：线性预测与 CEF}
\begin{homework}[练习 18]
\textsf{来源：《回归分析 60 题》已有原题第 3 页文件。}\par
设 $X,Y$ 的联合概率密度为 $f(x,y)=\frac32(x^2+y^2)$，其中 $0\leq x,y\leq1$（其余处为零）。求最优线性预测 $\alpha+\beta X$ 的系数，算出 $\mu(x)=\E[Y\mid X=x]$；最优线性预测与条件期望是否相同？
\end{homework}
\begin{solution}
积分得到
\[
\E X=\E Y=\frac58,\quad
\E X^2=\frac7{15},\quad
\E(XY)=\frac38.
\]
故
\[
\Var(X)=\frac{73}{960},\qquad
\Cov(X,Y)=-\frac1{64},
\]
\[
\beta=\frac{\Cov(X,Y)}{\Var(X)}=-\frac{15}{73},
\qquad
\alpha=\E Y-\beta\E X=\frac{55}{73}.
\]
最优线性预测为
\[
\widehat Y_L=\frac{55}{73}-\frac{15}{73}X.
\]
又
\[
f_X(x)=\frac32\left(x^2+\frac13\right),
\]
所以
\[
\mu(x)=\E(Y\mid X=x)
=\frac{x^2/2+1/4}{x^2+1/3}
=\frac{3(2x^2+1)}{4(3x^2+1)}.
\]
$\mu(x)$ 非线性，故条件期望与最优线性预测不同。
\end{solution}
\begin{dxtips}[线性预测与 CEF]\label{hw:reg:tip:18}
先用总体矩求仿射预测，再用条件密度求 CEF；两者优化的候选函数类不同。
\end{dxtips}

\subsection*{知识点定位}\begin{center}\begin{tabular}{ll}\toprule 题号 & 核心方法\\\midrule

8 & 回归误差的高阶矩 \\

9 & 乘法异方差模型 \\

10 & 遗漏交互项 \\

11 & 离散条件矩 \\

18 & 线性预测与 CEF \\

\bottomrule\end{tabular}\end{center}
\endgroup
\begingroup
\chapter*{第三章作业：古典线性回归模型}
\addcontentsline{toc}{chapter}{第三章作业：古典线性回归模型}
\markboth{第三章作业：古典线性回归模型}{第三章作业：古典线性回归模型}
\setcounter{section}{0}
\renewcommand{\thesection}{\arabic{section}}
\renewcommand{\thechapter}{3}
\renewcommand{\theHsection}{reg.homework.3.\arabic{section}}

\section{第三次作业}\label{hw:reg:assignment:3}
\setlength{\abovedisplayskip}{6pt}
\setlength{\belowdisplayskip}{6pt}
\setlength{\abovedisplayshortskip}{4pt}
\setlength{\belowdisplayshortskip}{4pt}

整理第二章剩余内容和第 3.2、3.3 节的全部已有阅读笔记。此次数学作业没有另列习题。总体矩公式要求有限二阶矩；总体 $R^2$ 要求 $\Var(Y)>0$；样本 $R^2$ 要求 TSS 正，OLS 与方差修正要求满列秩及 $n>K$。

\subsection*{第二章：条件期望函数的剩余内容}

\begin{proposition}[阅读笔记：条件方差、CEF 误差与全方差分解]\label{hw:reg:note:3:1}
设 $\mu(X)=\E(Y\mid X)$，定义 CEF 误差
\[
e=Y-\mu(X).
\]
则
\[
\E(e\mid X)=0,\qquad \E[h(X)e]=0
\]
对任意适当的可测函数 $h$ 都成立。条件方差为
\[
\sigma^2(X)=\Var(Y\mid X)=\E(e^2\mid X)
=\E(Y^2\mid X)-\mu(X)^2.
\]
总体波动可分解为
\[
\Var(Y)=\E\{\Var(Y\mid X)\}
+\Var\{\E(Y\mid X)\}.
\]
前者是平均条件噪声，后者是由解释变量引起的条件均值变动。
\end{proposition}

\begin{proposition}[阅读笔记：最优线性预测]\label{hw:reg:note:3:2}
当只允许使用线性函数 $X'b$ 预测 $Y$ 时，定义
\[
\beta=\arg\min_b\E[(Y-X'b)^2].
\]
目标函数对 $b$ 的一阶条件为
\[
\E[X(Y-X'\beta)]=0.
\]
若 $Q=\E(XX')$ 正定，则
\[
\beta=Q^{-1}\E(XY).
\]
记投影误差 $e_L=Y-X'\beta$，则
\[
\E(Xe_L)=0.
\]
这是一组无条件正交条件。若真实条件均值恰为线性函数，
\[
\E(Y\mid X)=X'\beta,
\]
则投影误差还满足更强的 $\E(e_L\mid X)=0$。
\end{proposition}

\begin{proposition}[阅读笔记：投影误差与拟合优度的总体形式]\label{hw:reg:note:3:3}
若 $X$ 含常数项，线性投影误差均值为零。由正交性，
\[
\E(Y^2)=\E[(X'\beta)^2]+\E(e_L^2).
\]
中心化后可写为
\[
\Var(Y)=\Var(X'\beta)+\Var(e_L).
\]
总体线性拟合优度可定义为
\[
R^2=\frac{\Var(X'\beta)}{\Var(Y)}
=1-\frac{\Var(e_L)}{\Var(Y)}.
\]
$R^2$ 描述线性投影解释的方差比例，但不能单独证明因果关系，也不能保证模型设定正确。
\end{proposition}

\begin{proposition}[阅读笔记：条件均值模型的正确设定]\label{hw:reg:note:3:4}
模型 $m(X,\theta)$ 正确设定，表示存在真值 $\theta_0$ 使
\[
\E(Y\mid X)=m(X,\theta_0).
\]
若错误地把非线性 CEF 设为线性函数，OLS 仍可估计最优线性近似，但系数只解释线性投影，不再等同于真实 CEF 的结构参数。判断模型设定时应：
\begin{enumerate}
  \item 根据理论选择解释变量和函数形式；
  \item 检查残差是否仍与 $X$ 的非线性变换相关；
  \item 比较不同设定的样本外预测与稳健性；
  \item 区分预测关系、相关关系和因果关系。
\end{enumerate}
\end{proposition}

\subsection*{第 3.2 节：普通最小二乘估计}

\begin{proposition}[阅读笔记：OLS 的矩阵推导]\label{hw:reg:note:3:5}
样本模型为
\[
Y=X\beta+u,
\]
其中 $Y$ 是 $n\times1$ 向量，$X$ 是满列秩的 $n\times K$ 矩阵。OLS 最小化残差平方和
\[
S(b)=(Y-Xb)'(Y-Xb).
\]
展开并求导：
\[
\frac{\partial S(b)}{\partial b}
=-2X'Y+2X'Xb.
\]
一阶条件给出正规方程
\[
X'X\widehat\beta=X'Y,
\]
因此
\[
\boxed{\widehat\beta=(X'X)^{-1}X'Y}.
\]
残差满足
\[
\widehat u=Y-X\widehat\beta,\qquad X'\widehat u=0.
\]
最后一个等式说明残差与每一列解释变量样本正交。
\end{proposition}

\begin{proposition}[阅读笔记：拟合值、残差和误差方差估计]\label{hw:reg:note:3:6}
定义
\[
P=X(X'X)^{-1}X',\qquad M=I-P.
\]
则
\[
\widehat Y=PY,\qquad \widehat u=MY.
\]
若经典同方差假设
\[
\E(u\mid X)=0,\qquad \Var(u\mid X)=\sigma^2I_n
\]
成立，则
\[
\widehat\sigma^2=\frac{\widehat u'\widehat u}{n-K}
\]
是 $\sigma^2$ 的无偏估计，因为
\[
\E(\widehat u'\widehat u\mid X)
=\E(u'Mu\mid X)
=\sigma^2\tr(M)
=\sigma^2(n-K).
\]
\end{proposition}

\subsection*{第 3.3 节：几何解释与拟合优度}

\begin{proposition}[阅读笔记：投影矩阵 $P$]\label{hw:reg:note:3:7}
$P$ 把任意 $Y\in\mathbb R^n$ 正交投影到 $X$ 的列空间：
\[
\widehat Y=PY\in\mathcal C(X).
\]
它具有
\[
P'=P,\qquad P^2=P,\qquad PX=X,
\]
且
\[
\operatorname{rank}(P)=\tr(P)=K.
\]
对称性表示正交投影，幂等性表示投影一次后再次投影不会改变结果。
\end{proposition}

\begin{proposition}[阅读笔记：消灭矩阵 $M$]\label{hw:reg:note:3:8}
$M=I-P$ 把 $Y$ 投影到 $\mathcal C(X)$ 的正交补：
\[
\widehat u=MY.
\]
它具有
\[
M'=M,\qquad M^2=M,\qquad MX=0,
\]
以及
\[
\operatorname{rank}(M)=\tr(M)=n-K.
\]
名称“消灭矩阵”来自 $MX=0$：它消去所有位于解释变量列空间中的分量。
\end{proposition}

\begin{proposition}[阅读笔记：几何分解与样本 $R^2$]\label{hw:reg:note:3:9}
OLS 给出正交分解
\[
Y=\widehat Y+\widehat u,\qquad
\widehat Y'\widehat u=0.
\]
若模型含截距，则 $\overline{\widehat Y}=\bar Y$、$\sum_i\widehat u_i=0$，并有
\[
\underbrace{\sum_i(Y_i-\bar Y)^2}_{TSS}
=
\underbrace{\sum_i(\widehat Y_i-\bar Y)^2}_{ESS}
+
\underbrace{\sum_i\widehat u_i^2}_{RSS}.
\]
因此
\[
R^2=\frac{ESS}{TSS}=1-\frac{RSS}{TSS}.
\]
加入解释变量不会增加 RSS，所以普通 $R^2$ 不会下降；比较不同维度模型时还应结合调整 $R^2$、理论合理性和样本外表现。
\end{proposition}

\begin{proposition}[阅读笔记：本次笔记的结构]\label{hw:reg:note:3:10}
\[
\text{条件期望}
\longrightarrow
\text{最优线性投影}
\longrightarrow
\text{样本 OLS}
\longrightarrow
\text{投影矩阵与 }R^2.
\]
第二章给出总体层面的最优预测与正交条件；第 3.2 节把总体问题变成样本最小二乘；第 3.3 节用投影几何解释拟合值、残差和拟合优度。
\end{proposition}


\endgroup

\end{document}
